% Genesis dynamics for the D-type dark sector: first-order confinement
% transition, chiral-reconnection baryogenesis, ADM co-genesis, two-epoch
% dark matter, and the open magnitude of eta (Paper XXII, DRAFT; companion
% to Paper XXI / dm/dtype_sector.tex)
%% hopfion-framework/preamble.tex
%% Shared preamble for all repository files.
%% \input this at the top of any standalone document.
%% Repository files are designed to be \input into papers directly.

\usepackage{amsmath,amssymb,amsthm}
\usepackage{hyperref}
\usepackage{enumitem}
\usepackage{booktabs}
\usepackage{xcolor}
\usepackage{geometry}
\geometry{margin=2.5cm}

%── Theorem environments ──────────────────────────────────────────────────────
\newtheorem{theorem}{Theorem}[section]
\newtheorem{lemma}[theorem]{Lemma}
\newtheorem{proposition}[theorem]{Proposition}
\newtheorem{corollary}[theorem]{Corollary}
\theoremstyle{definition}
\newtheorem{definition}[theorem]{Definition}
\newtheorem{result}[theorem]{Result}
\theoremstyle{remark}
\newtheorem{remark}[theorem]{Remark}
\newtheorem{observation}[theorem]{Observation}
\newtheorem{conjecture}{Conjecture}
\newtheorem{construction}[theorem]{Construction}
\newtheorem{condition}[theorem]{Condition}

%── Repository metadata commands ─────────────────────────────────────────────
% Usage: \status{published} or \status{open} or \status{conjectured}
\newcommand{\status}[1]{\par\noindent{\color{blue}\textbf{Status:} #1}}
% Usage: \source{Paper I, Theorem thm:scale}
\newcommand{\source}[1]{\par\noindent{\color{darkgray}\textbf{Source:} #1}}
% Usage: \residual{0.013\%} or \residual{exact}
\newcommand{\residual}[1]{\par\noindent{\color{teal}\textbf{Residual:} #1}}
% Usage: \depends{foundations/pentagon, foundations/suppression}
\newcommand{\depends}[1]{\par\noindent{\color{purple}\textbf{Depends on:} #1}}

%── Framework macros ─────────────────────────────────────────────────────────
\newcommand{\vp}{\varphi}
\newcommand{\bst}{\beta^{*}}
\newcommand{\Seff}{S_{\mathrm{eff}}}
\newcommand{\TCMB}{T_{\mathrm{CMB}}}
\newcommand{\Lcond}{\Lambda_{\mathrm{cond}}}
\newcommand{\LUV}{\Lambda_{\mathrm{UV}}}
\newcommand{\LQCD}{\Lambda_{\mathrm{QCD}}}
\newcommand{\MPl}{M_{\mathrm{Pl}}}
\newcommand{\ERy}{E_{\mathrm{Ry}}}
\newcommand{\aem}{\alpha_{\mathrm{em}}}
\newcommand{\sinW}{\sin^{2}\!\theta_{W}}
\newcommand{\vEW}{v_{\mathrm{EW}}}
\newcommand{\twoI}{2I}
\newcommand{\twoT}{2T}
\newcommand{\twoO}{2O}
\newcommand{\twoD}{2D}
\newcommand{\Ehat}{\hat{E}_8}
\newcommand{\QH}{Q_H}
\newcommand{\Ztwo}{\mathbb{Z}_2}
\newcommand{\Zthree}{\mathbb{Z}_3}
\newcommand{\SU}{\mathrm{SU}}
\newcommand{\SO}{\mathrm{SO}}
\newcommand{\cdir}{c_{\mathrm{dir}}}
\newcommand{\Odm}{\Omega_{\mathrm{DM}}}
\newcommand{\Ob}{\Omega_{b}}
\newcommand{\rCMB}{\rho_{\mathrm{CMB}}}
\newcommand{\Lobs}{\Lambda_{\mathrm{obs}}}
\newcommand{\rinfty}{\rho_{\infty}}
\newcommand{\oBD}{\omega_{\mathrm{BD}}}
\newcommand{\xNMC}{\xi_{\mathrm{NMC}}}
\newcommand{\Efb}{E_{\mathrm{fb}}}
\newcommand{\Kfb}{K_{\mathrm{fb}}}
\newcommand{\dimq}[1]{d_{#1}}
\newcommand{\deff}{d_{\mathrm{eff}}}
\newcommand{\Kth}{K_{\mathrm{th}}}
\newcommand{\Ja}{\mathcal{J}_a}
\newcommand{\Jfour}{\mathcal{J}_4}
\newcommand{\Jiso}{\mathcal{J}_{\mathrm{iso}}}
\newcommand{\rhoinf}{\rho_{\infty}}
\newcommand{\rhoa}{\rho_{\mathrm{atom}}}
\newcommand{\Heff}{H_{\mathrm{eff}}}
\newcommand{\epsr}{\varepsilon(r)}
\newcommand{\Hcond}{H_{\mathrm{cond}}}
\newcommand{\eps}{\varepsilon}
\newcommand{\Rzero}{R_0}
\newcommand{\Cstar}{C^*}
\newcommand{\mH}{m_H}
\newcommand{\aem}{\alpha_{\mathrm{em}}}
\newcommand{\ms}{\mu^*}
\newcommand{\lam}{\lambda}
\newcommand{\fzero}{f_0}
\newcommand{\ftwo}{f_2}
\newcommand{\Itor}[1]{I_{#1}^{\mathrm{tor}}}
\newcommand{\Jfb}{\mathcal{J}_{\mathrm{fb}}}
\newcommand{\bz}{\beta_0}
\newcommand{\sopt}{s_{\mathrm{opt}}}
\newcommand{\leff}{\lambda_{\mathrm{eff}}}
\newcommand{\ord}{\operatorname{ord}}
\newcommand{\G}{\mathcal{G}}
\newcommand{\Ia}{I_{2a}}
\newcommand{\Ib}{I_{4}}
\newcommand{\kk}{\kappa}
\newcommand{\aInv}{\alpha^{-1}}
\newcommand{\thetaG}{\theta_{\mathrm{g}}}
\newcommand{\LSC}{\mathrm{LSC}}
\newcommand{\mWZW}{m_{\mathrm{WZW}}}
\newcommand{\tang}{\dot\Gamma}
\newcommand{\Jones}{V}
\newcommand{\lk}{\mathrm{lk}}
\newcommand{\phistar}{\phi^*}
\newcommand{\CC}{\mathbb{C}}
\newcommand{\HH}{\mathbb{H}}
\newcommand{\OO}{\mathbb{O}}
\newcommand{\ZZ}{\mathbb{Z}}
\newcommand{\pcond}{p_{\mathrm{cond}}}
\newcommand{\pQM}{p_{\mathrm{QM}}}
\newcommand{\xcond}{x_{\mathrm{cond}}}
\newcommand{\CHSH}{\mathrm{CHSH}}
\newcommand{\Tsirelson}{2\sqrt{2}}
\newcommand{\RR}{\mathbb{R}}
\newcommand{\Hil}{\mathcal{H}}
\newcommand{\CQ}{\mathcal{C}_Q}
\newcommand{\CQr}{\mathcal{C}_Q^{\mathrm{red}}}
\newcommand{\bOmega}{\boldsymbol{\Omega}}
\newcommand{\dmu}{d\mu}
\newcommand{\nlvl}[1]{n\!\left(#1\right)}
\newcommand{\IE}{\mathrm{IE}}
\newcommand{\bsF}{{}_2F_1}
\newcommand{\Qgr}{Q_{\mathrm{group}}}
\newcommand{\Mup}{M_{\mathrm{up}}}
\newcommand{\Mdn}{M_{\mathrm{down}}}
\newcommand{\Tw}{\mathrm{Tw}}
\newcommand{\alphas}{\alpha_s}
\newcommand{\alphaem}{\alpha_{\mathrm{em}}}




The deconfinement free energy is a Polyakov-loop Landau--Ginzburg
potential for the order parameter $L$, constrained by a $\Zthree$
center symmetry~\cite{Pisarski2000}:
\begin{equation}
  V(L)=\tfrac{b_2(T)}{2}|L|^2-\tfrac{b_3}{3}(L^3+L^{*3})+\tfrac{b_4}{4}|L|^4-h(L+L^*).
  \label{P22:eq:polyakov}
\end{equation}
The cubic $b_3$ (allowed only by $\Zthree$) drives a first-order
transition; the explicit breaking $h$ weakens it, and above a critical
$h_c$ the transition becomes a crossover (the physical light-quark
point of real QCD).

\begin{proposition}[Framework first-order]
\label{P22:prop:fo}
In the density-feedback Hopfion the $\Zthree$ is the colour symmetry of
the trefoil $T_{2,3}$ (the $\Zthree$ outer automorphism of $\twoT$,
Paper~XV), so the cubic $b_3\neq0$ is present by construction (the
colour-singlet baryon is the $\Zthree$ invariant
$3\times3\times3\to1$). The explicit breaking is small,
$h\sim e^{-E_{\rm sb}/T}$ with $E_{\rm sb}\sim2m_{\rm const}\sim600$~MeV
the string-breaking energy and $T\sim150$~MeV, giving
$h\sim0.02$--$0.14$, below the critical endpoint $h_c\approx0.21$. The
transition is therefore first-order --- the framework differs from
real QCD (crossover) precisely in having no free light quarks:
$\Zthree$ is broken only by string-breaking, not by light dynamical
quarks. A density-feedback contribution (the $(1+\beta_*\rho)$
softening) adds to the coefficients and provides a second, independent
first-order driver (\path{src_paper22/firstorder_transition.py}); the
first-order region is an intermediate window in $\beta_*\rho$
(P22:OP1).
\end{proposition}
\status{draft}
\source{P22:prop:fo}

The coefficients of equation~\eqref{P22:eq:polyakov} are not free: they
are the holonomy expansion of the thermal free energy of the condensate
excitations, which carry the \emph{semi-Dirac} dispersion of the medium
(Paper~VII),
\begin{equation}
  E(k,\theta)^2=A^2k^4\sin^4\theta+v^2k^2\cos^2\theta ,
  \label{P22:eq:semidirac}
\end{equation}
linear along the easy (radial) axis and quadratic along the hard
(perpendicular) axis. The anisotropic phase space changes the
temperature power: $\int\!d^3k\,e^{-nE/T}\propto(T/n)^{2}$ rather than
$(T/n)^3$, so the single-phase holonomy potential is
\begin{equation}
  f(T,\chi)=-\sum_{n\ge1}c_n\cos n\chi,\qquad c_n\propto\frac{T^3}{n^{3}} ,
  \label{P22:eq:weiss}
\end{equation}
a cubic-Bernoulli ($B_3$) Weiss potential, in contrast to the
relativistic quartic ($B_4$, $c_n\propto T^4/n^4$). Summing the three
$\SU(3)$ colour (holonomy) phases $(\chi,0,-\chi)$ along the real axis,
with $L=(1+2\cos\chi)/3$ and $\cos n\chi=T_n\!\big(\tfrac{3L-1}{2}\big)$
(Chebyshev), and expanding in $L$ gives
\begin{equation}
  \frac{b_3}{b_4}\approx0.17 \quad(\text{semi-Dirac }B_3),\qquad\text{vs}\quad 0.16\ (B_4),
  \label{P22:eq:b3b4}
\end{equation}
so the epoch's geometry modestly \emph{enhances} the cubic (strengthens
first-order) relative to a relativistic gas. Confinement gaps the
excitations through a $\Zthree$-even coupling to the order parameter,
$\Delta(L)^2=\Delta_0^2(1-L^2)$ with $\Delta_0\sim\Lcond(\mathrm{QCD})$;
this contributes a thermal mass $b_2(T)$ whose sign changes at
$T_c\sim2\Lcond\sim T_{\rm QCD}$, matching the framework's near-critical
condition $T/\Lcond\approx2$ (Paper~VII). (Computed in
\path{src_paper22/semidirac_thermal.py},
\path{src_paper22/su3_colorsum_coeffs.py}, and
\path{src_paper22/gap_b2T.py}.)
\status{draft}
\source{P22:subsec:coeffs}

\begin{theorem}[Derived first-order transition]
\label{P22:thm:fo}
The full thermal free energy combining the holonomy and the gap in one
normalisation,
\begin{equation}
  F(T,L)=-T\sum_{n\ge1}\frac1n\,I_n\!\big(T,\Delta(L)\big)\,[\,2\cos n\chi+1\,]+a_{\rm vac}L^2,
  \qquad I_n=\int\!\frac{d^3k}{(2\pi)^3}\,e^{-nE/T},
  \label{P22:eq:FTL}
\end{equation}
has, below $T_c$, its global minimum on the confined branch
($L\!\approx\!0.35$--$0.45$) and, above $T_c$, at the deconfined value
$L\!\to\!1$, with the two degenerate at $T_c\sim(0.9$--$1.5)\,\Delta_0
\sim O(\Lcond)$. The order parameter therefore jumps
\emph{discontinuously} by $\Delta L\approx0.55$--$0.65$ across a finite
barrier: the transition is first-order, driven by the $\Zthree$ cubic
(the $n=3k$ holonomy harmonics, i.e.\ the trefoil's three-fold colour
structure). (Computed in \path{src_paper22/unified_FTL.py}.) The jump
in the order parameter is large (strongly first-order in
order-parameter space), while the latent heat is a modest fraction of
the radiation bath, $\alpha=L_{\rm heat}/\rho_{\rm rad}\approx0.03$--$0.06$,
giving bubble-wall speeds $v_w\sim0.15$--$0.2$ (deflagration).
\end{theorem}
\status{draft}
\source{P22:thm:fo}

A baryon is a chiral wound state: the trefoil (writhe $+3$) differs
from its mirror (Paper~XVIII). The three Sakharov conditions map onto
knot dynamics: baryon-number violation is the reconnection
$\QH{=}2+\QH{=}2\to\QH{=}3+\QH{=}1$ (baryons form from leptons by
threading); C and CP violation are the derived Wess--Zumino--Witten
phase $\delta_{\rm CP}=65.0^\circ$ (Paper~V) together with the chiral
writhe; and the departure from equilibrium is the first-order
transition of Theorem~\ref{P22:thm:fo}. Because the reconnection rate
is a strong-interaction rate, $\Gamma/H\sim10^{19}$ at $T_{\rm QCD}$
(deep equilibrium), only a genuine first-order freeze-out escapes the
resulting washout. Here that escape is structural: a reconnection needs
free trefoil strands, which exist only in the deconfined phase, so once
the bubble of confined phase forms the baryon-number-violating rate
shuts off \emph{automatically} --- the analogue of the electroweak
``sphaleron shutoff in the broken phase'' is guaranteed by confinement
itself, with no condition on the transition strength. Since
$\sin\delta_{\rm CP}=0.91$ is of order unity, the CP bias is not the
limiting factor.
\status{draft}
\source{P22:sec:baryo}

The phase $\delta_{\rm CP}$ is a property of the shared $\twoI$ vacuum,
which contains both $\twoT$ (baryon) and the D-type $\twoD_n$ (dark) as
subgroups. The same chiral, CP-violating reconnection at the bubble
walls therefore biases \emph{both} freeze-outs, seeding correlated
asymmetries. The D-type sector is Dirac (its odd-rank spinor is
complex, Paper~XXI), so it carries a dark number and an asymmetry; its
symmetric part self-annihilates (the sector is self-interacting through
the shared director field), leaving the asymmetric residue. Its
abundance follows the asymmetric-dark-matter arithmetic,
\begin{equation}
  \frac{\Odm}{\Ob}=r\,\frac{m_{\rm dark}}{m_p},\qquad r=\frac{n_{\rm dark}}{n_b},
  \label{P22:eq:adm}
\end{equation}
with $r$ the ratio of the dark and baryon number asymmetries. One
feature is robust: $r$ is \emph{independent of the unpinned magnitude}
of $\eta$ (P22:OP1), because both asymmetries are seeded by the same
$\delta_{\rm CP}$ at the same transition, so the common transport
efficiency cancels in the ratio. What is \emph{not} derived is the
value of either $r$ or $m_{\rm dark}$: $r\sim O(1)$ is taken on the
ground that $\twoD_3$ (order $12$) and $\twoT$ (order $24$) are
comparable subgroups of $\twoI$, an estimate from subgroup orders, not
a computed branching; and because the D-type carries no
$Q_{\rm group}$ (Paper~XXI), its mass tower is undefined and the
framework does \emph{not} pin $m_{\rm dark}$. With $r\sim O(1)$,
matching the observed $\Odm/\Ob\approx5.4$ then \emph{fixes}
$m_{\rm dark}\approx5$~GeV --- the canonical asymmetric-dark-matter
scale, a sterile $\sim$few-GeV state that evades direct detection by
construction. This is a \emph{conditional consistency}, not a
parameter-free prediction: it rests on $r\sim O(1)$, on efficient
symmetric self-annihilation through the sector's SIDM channel, and on
the dark relic being stable (the existence question of Paper~XXI).
\status{draft (conditional consistency, not a parameter-free prediction)}
\source{P22:eq:adm, P22:sec:abund}

The stable D-type relic and the baryon-sourced director \emph{strain}
of Paper~VII are two configurations of the one condensate, dominating
the dark-matter phenomenology at two different epochs --- a single
fabric in two modes, not two dark species. Paper~VII identifies the
\emph{galactic} dark matter as the director strain, but that component
is late and isocurvature: its ultralight scale $m_\xi\sim H_0$ keeps it
Hubble-frozen until $z\sim0$, so it cannot supply the cold matter the
acoustic peaks require at recombination. The D-type relic is precisely
the missing ingredient --- non-baryonic, electromagnetically and weakly
sterile (Paper~XXI), non-ultralight ($m_{\rm dark}\sim5$~GeV
$\gg H(z_{\rm rec})$), cold, and, being seeded by the co-genesis above
correlated with the baryon asymmetry, \emph{adiabatic} rather than
isocurvature: it clusters as standard cold dark matter at
recombination and, at the cosmic ratio $\Odm/\Ob\approx5$, supplies the
cluster lensing mass where the baryon-tracking strain alone falls
short. At $z\sim0$ the strain unfreezes and yields the tight
radial-acceleration relation (RAR) because it tracks the baryons
(Paper~VII). A cosmic-ratio relic, were it collisionless here too,
would build halos that do not track baryons and would spoil that
tightness; it does not, because the relic self-interacts through the
\emph{shared} director field --- the single $\hat n$ that carries the
Bell correlations of Papers~VIII--IX --- and that interaction inherits
the field's anisotropy via the same kernel
$S_{\rm eff}=\sin^4\theta/[\vp^6(1+\beta_*\rho)]$ that gates those
correlations: resolved in a coherent geometry (inside a galaxy, where
it relaxes the relic to the baryon-dominated potential, preserving the
RAR, with weak coupling $C\sim10$--$10^3$~eV and dark sound speed
$c_s^2=C/m_{\rm dark}\sim10^{-8}$), washed out under isotropic averaging
(the smooth recombination medium and the inter-galactic field, where
the relic is inert and collisionless, with no cosmological
over-closure). This reproduces the observed division ---
$\Lambda$CDM-like at the microwave background and in clusters, a
MOND-like RAR in galaxies --- from two modes of one condensate, and
ties the dark-matter self-interaction to the same geometry-dependent
$\sin^4\theta$ structure as the Bell correlations: one anisotropy, two
arenas. It is conditional on the relic's existence and abundance
(Paper~XXI). The residual open question is the \emph{within-halo}
energetics: whether the cumulative interaction energy is simply the
already-counted coherent halo strain plus a finite perturbative
relaxation, or requires a further nonlinear ($S^2$) saturation.
Falsifiable signatures shared across domains: the
planar-versus-spherical Bell geometry-dependence (testable with 3D or
satellite entanglement protocols) probes the same kernel that gates the
dark sector; a sterile $\sim$few-GeV self-interacting relic; a mild
suppression of sub-megaparsec power; and weak lensing exceeding the
kinematic mass in galaxies by $O((c_{\rm dir}/v)^2)$.
\status{draft (conditional on the relic's existence and abundance)}
\source{P22:sec:twoepoch}

All three Sakharov conditions are derived: the first-order transition
from the epoch's thermal free energy (Theorem~\ref{P22:thm:fo}), the
order-unity CP phase $\delta_{\rm CP}=65.0^\circ$, and the automatic
shutoff of baryon-number violation in the confined phase. What is
\emph{not} derived is the \emph{magnitude} of the baryon asymmetry
$\eta=7.04\,n_B/s$. Writing
\begin{equation}
  \frac{n_B}{s}\sim A_{\rm pref}\,\varepsilon_{\rm CP}\,\delta_{\rm neq}\,b_{\rm ch},\qquad
  A_{\rm pref}=\frac{15}{4\pi^2 g_*}\sim2\times10^{-2},\qquad
  \varepsilon_{\rm CP}=\sin\delta_{\rm CP}\,c_{\rm int},
  \label{P22:eq:eta}
\end{equation}
two factors are \emph{not} small: the departure from equilibrium at the
wall, $\delta_{\rm neq}=v_w/(\Gamma L_w)\sim10^{-2}$ (set by the derived
$v_w\sim0.2$), and the CP interference, $c_{\rm int}\sim0.6$ (the
argument of a two-amplitude sum, Paper~V). With a naive isotropic
channel fraction $b_{\rm ch}\sim0.2$, these factors \emph{overproduce}
$\eta$ by $\sim10^{4}$--$10^{5}$: reconnection baryogenesis is
efficient, so $\eta$'s problem is a \emph{suppression}, not a
smallness. (\path{src_paper22/transport_eta.py},
\path{src_paper22/branching_eta.py}.)

The suppression is supplied by the \emph{anisotropy}. The baryon
($(E_6)_1$, $Q_H{=}3$) is the \emph{heavy band} of the semi-Dirac
dispersion~\eqref{P22:eq:semidirac}, carrying the directional
suppression of the condensate (Paper~VII, Axiom~A2),
\begin{equation}
  S_{\rm eff}=\frac{\sin^4\theta}{\vp^6\,(1+\beta_*\rho)},
  \label{P22:eq:seff}
\end{equation}
whereas the competing light-band channels do not. The baryon-forming
reconnection therefore picks up $b_{\rm ch}\to b_{\rm ch}\,S_{\rm eff}^{\,p}$,
with $p$ the number of heavy-band projections in the final state; the
factor $\vp^6$ is the Wess--Zumino--Witten/Bogomolny power (physical,
not frame-dependent), giving a real suppression of
$\sim\vp^{-6}\approx0.056$ per power. Crucially, $\beta_*\rho$ in
equation~\eqref{P22:eq:seff} is the \emph{same} density-feedback
combination that drives the first-order transition (the
$(1+\beta_*\rho)$ softening of Proposition~\ref{P22:prop:fo}): one
variable doing two jobs. Because $\rho$ tracks $T^4$, $\beta_*\rho$ is
large at the dense QCD epoch --- the source of the baryon-channel
suppression --- but cannot be arbitrarily large: deep saturation
softens the barrier away and returns a crossover, so a first-order
transition exists only in an \emph{intermediate window} of
$\beta_*\rho$. Overlaying the two constraints: for $p=2$ a
self-consistent value ($\beta_*\rho\sim15$) lies \emph{inside} the
first-order window and gives $\eta\approx6\times10^{-10}$, with the
window bracketing the observed value; for $p=1$ the observed $\eta$
falls at the over-saturated edge of the window, a mild overproduction.
A self-consistent solution thus exists but is not forced.
(\path{src_paper22/branching_anisotropic.py},
\path{src_paper22/beta_rho_window.py}.) The magnitude of $\eta$ has
thus reduced from a transport unknown spanning many orders to a single
density-feedback quantity, $\beta_*\rho$ at the QCD epoch, co-constrained
by the requirement of a first-order transition and shared with the
density screening of the dark sector. What remains genuinely open is
(i) the \emph{location} of the first-order window in physical units
(the ratio $\gamma/b$ of the feedback free energy, not yet fixed) and
(ii) the heavy-band power $p$ (P22:OP1). The co-genesis ratio
$\Odm/\Ob\sim O(5)$ is independent of all of this and stands as the
sector's falsifiable prediction.
\status{draft (open: $\eta$'s magnitude pinned only up to the window location and power $p$)}
\source{P22:sec:open, P22:eq:eta, P22:eq:seff}
